% !PRINTED 672
\weberpage{672}
%<<CONT
2. \gesperrt{Sind zwei $n$-Ecke mit einem dritten äquivalent, so sind sie auch untereinander äquivalent. Man vereinigt danach äquivalente Polygone $\mathfrak{A}, \mathfrak{A}', \ldots$ in einer Polygonklasse $A$. Jedes Polygon ist in einer, und nur in einer Klasse enthalten.}

Es existieren auch Polygone, die mit keinem anderen äquivalent sind, und die daher jedes für sich eine besondere Klasse bilden. Diese heißen \gesperrt{isolierte Polygone}.

3. \gesperrt{Ist $\mathfrak{A} \sim \mathfrak{A}'$, $\mathfrak{B} \sim \mathfrak{B}'$, so ist $\mathfrak{A}\mathfrak{B} \sim \mathfrak{A}'\mathfrak{B}'$.}

Dies folgt aus

\weq{4}{\eta = \dfrac{\mathfrak{A}}{\mathfrak{A}'}\,, \quad \eta_1 = \dfrac{\mathfrak{B}}{\mathfrak{B}'}\,, \quad \eta\eta_1 = \dfrac{\mathfrak{A}\mathfrak{B}}{\mathfrak{A}'\mathfrak{B}'}.}

Die Klasse $C$, der das Produkt $\mathfrak{A}\mathfrak{B}$ aus den Klassen $A, B$ angehört, enthält daher alle Produkte aus einem Polygon der Klasse $A$ mit einem Polygon der Klasse $B$ (unter Umständen auch noch andere) und ist daher durch die Klassen $A$ und $B$ vollständig bestimmt. Darauf beruht die Multiplikation (Komposition) der Klassen, die sich in der Formel ausdrückt:

\weq{5}{C = AB = BA.}

Aus einer Gleichung zwischen drei Klassen $A, B, M$

\[
MA = MB
\]

folgt hiernach $A = B$. Denn ist $\mathfrak{M}\mathfrak{A} \sim \mathfrak{M}\mathfrak{B}$, so folgt aus (2) $\mathfrak{A} \sim \mathfrak{B}$. Ist $\mathfrak{A} \sim \mathfrak{A}'$ und

\[
\mathfrak{A}\mathfrak{B} = \mathfrak{C} \quad \text{und} \quad \mathfrak{A}'\mathfrak{B} = \mathfrak{C}'\,,
\]

so ist auch $\mathfrak{C} \sim \mathfrak{C}'$. Wenn also $\mathfrak{A}$ in $\mathfrak{C}$ enthalten ist, so ist jedes mit $\mathfrak{A}$ äquivalente Polygon in einem mit $\mathfrak{C}$ äquivalenten Polygon $\mathfrak{C}'$ enthalten.

4. \gesperrt{Wir nennen eine Klasse $C$ durch eine Klasse $A$ teilbar, wenn ein Polygon von $A$ in einem Polygon von $C$ enthalten ist, und setzen}

\weq{6}{AB = C.}

\gesperrt{Die Klasse $B$ ist durch $A$ und $C$ vollständig bestimmt.}

\begin{center}
\textbf{§ 192. Polygonscharen.}
\end{center}

Wenn $\mathfrak{A}, \mathfrak{A}_1, \mathfrak{A}_2, \ldots, \mathfrak{A}_s$ Polygone einer Klasse $A$ sind, so gibt es $s$ Funktionen in $\Omega$:

\weq{1}{\eta_1 = \dfrac{\mathfrak{A}_1}{\mathfrak{A}}\,, \quad \eta_2 = \dfrac{\mathfrak{A}_2}{\mathfrak{A}}\,, \quad \cdots, \quad \eta_s = \dfrac{\mathfrak{A}_s}{\mathfrak{A}}\,.}
